\documentclass[12pt]{article}
\usepackage[utf8]{inputenc}
\usepackage[T1]{fontenc}
\usepackage{hyperref}
\title{Bunganditj Language Resources for K-12 Curriculum Development}
\author{Compiled by Claude Sonnet 5.5 (claude-sonnet-5-5) for Andr\'as Kornai}
\date{}
\begin{document}
\maketitle

\section{Introduction}
Bunganditj (also Buwandik, Boandik, Bungandidj; ISO 639-3 xbg) is a Pama-Nyungan language of the Mount Gambier region of South Australia and western Victoria. It has no living native speakers and no recordings of speakers. A community revival and maintenance programme is under way.

\section{Primary Educational Resources}
No school primers or textbooks were found. The Bunganditj Language Reclamation Working Party produced a 2012 book, \emph{Boandik/Bunganditj language reclamation project}, which includes word lists. I saw no URL for it in search results.

\section{Linguistic and Reference Materials}
\begin{itemize}
\item Barry Blake (2003), \emph{The Bunganditj (Buwandik) language of the Mount Gambier Region}, Pacific Linguistics: word lists, pronunciation guide and English translation. It is catalogued at \url{https://phoible.org/sources/blake_bunganditj_2003}.
\item Wikipedia, ``Bungandidj language''. \url{https://en.wikipedia.org/wiki/Bungandidj_language}
\item Mary-Anne Gale, a paper on using archival material to build new vocabulary (University of Adelaide AUSTRALEX). \url{https://www.adelaide.edu.au/australex/publications/gale.pdf}
\item ASJP language page for Bunganditj. \url{https://asjp.clld.org/languages/BUNGANDITJ}
\end{itemize}
Blake's separate ``Grammar and glossary of Buwandik/Bunganditj'' manuscript is held by ASEDA. I saw no URL for it. The historical sources are short sketches by D. S. Stewart and R. H. Mathews and about twelve vocabulary lists.

\section{Religious and Community Resources}
The Burrandies Aboriginal Corporation and the Bunganditj Language Reclamation Committee, made up of Boandik community members, lead the revival work. The Mobile Language Team has a Bunganditj/Boandik page. \url{https://mobilelanguageteam.com.au/languages/bunganditj-boandik/}

\section{Digital and Audio Resources}
No dictionary app, online dictionary or audio was found. In 2022 Mount Gambier City Council adopted a Dual Naming Policy, and six local sites received Bunganditj names alongside their English ones, including Blue Lake (Warwar) and Cave Garden (Thugi). \url{https://theleadsouthaustralia.com.au/?p=41832}

\section{Curriculum Development Context}
The language is reconstructed from archival sources, and the community has led the revival for over a decade. The dual-naming sites could be used in place-based lessons. Any curriculum should be developed with Burrandies Aboriginal Corporation and the Reclamation Committee.
\end{document}